% VI24_FULL_SOLUTION_REFINEMENT_V1
\subsection*{VI.24.P15: Nilpotents under field extension}
\begin{problem}\label{prob:vi24-p15}
Show that dual numbers remain nonreduced after any field extension $k\subseteq K$.
\end{problem}

\ifdefined\FullProblemDossiers
\paragraph{Why this problem matters.}
This dossier isolates a reusable piece of the base-change calculus and records both its categorical and affine content.

\paragraph{Useful ideas before starting.}
Start from the Cartesian universal property.  When the schemes are affine, translate the square contravariantly into tensor products, quotients, or localizations, and then translate the resulting ring statement back into geometry.

\paragraph{Guided hints.}
\begin{enumerate}
\item Identify the relevant pullback or scalar-extension ring.
\item Use the appropriate universal property or exactness statement.
\item Check that the canonical maps are the expected projections.
\item State the geometric consequence, not only the ring isomorphism.
\end{enumerate}

\begin{solution}
Let $B=k[\varepsilon]/(\varepsilon^2)$.  By scalar extension of a polynomial quotient,
\[
B\otimes_kK\cong K[\varepsilon]/(\varepsilon^2).
\]
The class $\bar\varepsilon$ is nonzero: the quotient has $K$-basis $1,\bar\varepsilon$.  But
$\bar\varepsilon^2=0$.  Hence the scalar extension is nonreduced.  In scheme language, the doubled point
\[
\Spec k[\varepsilon]/(\varepsilon^2)
\]
remains a doubled point after every field extension.  This example should be contrasted with P16: P15 starts nonreduced and stays nonreduced, whereas P16 starts with reduced fields and creates nilpotents through a purely inseparable tensor product.
\end{solution}

\paragraph{What happened mathematically.}
The base-changed object is forced by a universal property; tensor products make that universal property computable on affine charts.

\paragraph{Extensions.}
The same pattern will be reused in VI/25 for fibers and in VI/26 for properties of morphisms under change of base.

\paragraph{Provenance.}
Canonical persistence-of-nilpotents example.
\else
\paragraph{Provenance.} Canonical persistence-of-nilpotents example.
\fi
